\section*{Chapitre \ref{ch:b3:genetic-engineering} --- Génie génétique et biotechnologies}

\begin{solution}{exo:b3:genetic-engineering:1}
Une origine de réplication, pour que l'hôte recopie le \omterm{def:b3:genetic-engineering:tools}{plasmide} et le
transmette à ses cellules filles~; un \omterm{def:b3:genetic-engineering:tools}{marqueur de sélection},
d'ordinaire une résistance à un antibiotique, pour que seules les
cellules porteuses du \omterm{def:b3:genetic-engineering:tools}{plasmide} poussent~; et un site de restriction
unique (ou un groupe de sites) où l'insert est ligaturé sans couper le
\omterm{def:b3:genetic-engineering:tools}{plasmide} ailleurs.
\end{solution}

\begin{solution}{exo:b3:genetic-engineering:2}
Une séquence donnée de six bases apparaît une fois tous les
$4^{6} = 4096$ bp dans un ADN aléatoire~: environ
$4.6\times 10^{6}/4096 \approx 1100$ fragments de \qty{4.1}{kb} en
moyenne. Les génomes réels ne sont pas aléatoires --- la composition en
bases, l'usage des codons et l'évitement de certaines séquences
déplacent le compte (le nombre réel de sites EcoRI chez \emph{E.~coli}
est de quelques centaines).
\end{solution}

\begin{solution}{exo:b3:genetic-engineering:3}
La dénaturation à \qty{95}{\celsius} sépare les brins~; l'appariement à
\qtyrange{50}{65}{\celsius} laisse les \omterm{def:b3:genetic-engineering:pcr}{amorces} s'apparier à leurs
sites~; l'élongation à \qty{72}{\celsius} laisse la polymérase copier
depuis chaque \omterm{def:b3:genetic-engineering:pcr}{amorce}. Une polymérase ordinaire serait détruite à
\qty{95}{\celsius} dès le premier cycle et devrait être ajoutée à neuf
trente fois~; la Taq survit à tous les cycles.
\end{solution}

\begin{solution}{exo:b3:genetic-engineering:4}
Un \omterm{def:b3:genetic-engineering:crispr}{ARN guide} dont les vingt bases correspondent à la cible, et un \omterm{def:b3:genetic-engineering:crispr}{PAM}
(NGG) immédiatement en 3$'$ de la cible sur le brin non ciblé. Après la
\omterm{def:b3:dna-repair:lesions}{cassure double brin}~: la jonction d'extrémités, qui laisse une petite
insertion ou délétion et détruit en général le cadre de lecture du
gène~; ou la \omterm{def:b3:dna-repair:dsb}{recombinaison homologue} avec une matrice portant la
séquence désirée, qui l'y écrit.
\end{solution}

\begin{solution}{exo:b3:genetic-engineering:5}
$50\times 1.9^{35} = 50\times 5.7\times 10^{9} \approx 3\times 10^{11}$
copies. $\Delta C_{T} = 7.3$~; rapport $1.95^{7.3} = e^{7.3\ln 1.95}
\approx 130$~: le premier échantillon avait environ $130$ fois plus de
matrice.
\end{solution}

\begin{solution}{exo:b3:genetic-engineering:6}
Les bactéries ne savent pas épisser~: une copie génomique avec ses
introns serait transcrite en un message inutilisable, on emploie donc
l'ADNc sans intron, recopié sur l'ARNm mature. Autres obstacles~: le
promoteur humain n'est pas reconnu par la polymérase bactérienne
(fournir un promoteur bactérien et un site de fixation du ribosome)~;
l'usage des codons humains diffère (synthétiser le gène avec les codons
préférés de l'hôte)~; la protéine peut ne pas se replier ou s'agréger
(abaisser la température, la fusionner à un partenaire soluble, ou la
replier depuis les corps d'inclusion)~; la glycosylation manque
(employer la levure ou des cellules de mammifère si elle compte)~; les
ponts disulfure exigent le périplasme oxydant.
\end{solution}

\begin{solution}{exo:b3:genetic-engineering:7}
Les cellules souches ciblées sont injectées dans un blastocyste hôte et
se mêlent à ses propres cellules, de sorte que la souris est bâtie de
deux génotypes --- une chimère --- et elle ne transmet l'allèle que si
une partie de ses cellules germinales dérive des cellules ciblées.
Croiser la chimère avec une souris sauvage donne des hétérozygotes (à
partir de la moitié de ses gamètes si la lignée germinale est à moitié
ciblée)~; l'intercroisement d'hétérozygotes donne un quart
d'homozygotes \omterm{def:b3:genetic-engineering:transgenic}{knock-out}.
\end{solution}

\begin{solution}{exo:b3:genetic-engineering:8}
Exact~: $6.4\times 10^{9}\times 4^{-20}\times (1/16) \approx 4\times
10^{-4}$ site (le GG du \omterm{def:b3:genetic-engineering:crispr}{PAM} coûte un facteur $16$) --- aucun, en
pratique. À deux mésappariements près il y a $1 + 60 + 1710 = 1771$
séquences~: $1771\times 4\times 10^{-4} \approx 0.7$ site fortuit. Le
génome réel n'est pas aléatoire, et l'on vérifie un guide contre lui
directement.
\end{solution}

\begin{solution}{exo:b3:genetic-engineering:9}
$p_{1} = 0.05 + 0.8\times 0.05\times 0.95 = 0.088$~; $p_{2} = 0.152$~;
$p_{3} = 0.255$~; puis $0.41$, $0.60$~: environ cinq générations pour
dépasser $0.5$, contre l'estimation
$\ln(1/(2\times 0.05))/\ln 1.8 \approx 4$.
\end{solution}

\begin{solution}{exo:b3:genetic-engineering:10}
La jonction d'extrémités agit dans toute cellule, y compris les
cellules souches quiescentes où la \omterm{def:b3:dna-repair:dsb}{recombinaison homologue}, propre aux
phases S et G2, est inefficace~; elle n'exige la livraison d'aucune
matrice et son produit --- n'importe quelle interruption de
l'amplificateur --- fait l'affaire, alors que la correction exige la
séquence exacte et ne réussit que dans une minorité de cellules. De
plus, la même édition vaut pour toute mutation de la $\beta$-globine,
drépanocytaire ou thalassémique. Inconvénient~: elle supprime un
élément régulateur (avec tous les autres rôles qu'il peut avoir), elle
laisse l'allèle drépanocytaire en place, et le mélange d'insertions et
de délétions n'est pas maîtrisé.
\end{solution}

\begin{solution}{exo:b3:genetic-engineering:11}
Avec une valeur sélective $1 - s$ chez l'homozygote~: $p' = \bigl[p^{2}(1-s) + p(1-p)(1+c)
\bigr]/(1 - sp^{2})$. Depuis la rareté les homozygotes sont
négligeables et le forçage croît comme $(1+c)p$ quel que soit $s$~; sa
fixation se décide près de $p = 1$, où un allèle sauvage rare de
fréquence $q$ revient comme $q' \approx q(1-c)/(1-s)$~: fixation si
$c > s$, équilibre intermédiaire sinon. Résistance~: chaque conversion
manquée (jonction d'extrémités au lieu de recombinaison) laisse un
allèle que le guide ne reconnaît plus~; si de tels allèles ne coûtent
rien --- une insertion ou délétion en phase dans une région non
essentielle --- ils n'ont aucun coût là où le forçage en a un, de sorte
que la sélection les favorise et qu'ils le remplacent. Les remèdes sont
de viser des séquences essentielles où toute insertion ou délétion est
létale, et d'employer plusieurs guides à la fois.
\end{solution}

\begin{solution}{exo:b3:genetic-engineering:12}
Une \omterm{def:b3:genetic-engineering:crispr}{Cas9} injectée dans un zygote peut agir après les premières
divisions, de sorte que des cellules différentes reçoivent des
réparations différentes --- le mosaïcisme~; la réparation de chaque
cassure est choisie par la cellule, la jonction d'extrémités donnant
des insertions et délétions imprévisibles, de grandes délétions ou une
perte d'hétérozygotie au point de coupure~; l'issue par \omterm{def:b3:dna-repair:dsb}{recombinaison
homologue} reste minoritaire~; des cassures hors cible surviennent en
des sites qu'on ne peut pas tous prédire~; et l'embryon ne peut être
vérifié qu'en séquençant quelques cellules biopsiées, qui ne parlent
pas pour les autres. Preuves nécessaires~: des méthodes qui éditent
toutes les cellules à l'identique (avant la première phase~S) avec des
rendements proches de \qty{100}{\%}, le séquençage du génome entier de
chaque cellule d'embryons animaux édités montrant l'absence de tout
changement non voulu, et le suivi à long terme d'animaux édités sur
plusieurs générations --- dont rien n'existe.
\end{solution}

\begin{solution}{pb:b3:genetic-engineering:1}
\textbf{1.} Une fois tous les $4096$ bp. $P(\text{aucun site}) = (1 - 1/4096)^{1500}
= e^{-0.37} = 0.69$.
\textbf{2.} Choisir une enzyme dont le site est absent du gène~; ou
amplifier le gène par PCR avec des \omterm{def:b3:genetic-engineering:pcr}{amorces} portant de nouveaux sites de
restriction à leur extrémité 5$'$ (ou recourir au clonage à extrémités
franches ou indépendant de la ligature).
\textbf{3.} \qty{0.1}{\micro g}\,$\times 2\times 10^{7} = 2\times
10^{6}$ colonies~; \qty{10}{\%}, soit $2\times 10^{5}$, portent
l'insert.
\textbf{4.} Blanches~: \qty{10}{\%}. La moitié des inserts est dans le
bon sens, donc $P(\text{aucun parmi } n \text{ blanches}) = 0.5^{n} \le 0.01$
donne $n \ge 6.6$~: en prélever sept.
\textbf{5.} $32$ doublements, $2^{32} = 4.3\times 10^{9}$ cellules,
$2.1\times
10^{11}$ \omterm{def:b3:genetic-engineering:tools}{plasmides}~; chacun $5500\times 650 = 3.6\times 10^{6}$ Da $=
5.9\times 10^{-18}$ g~: \qty{1.3}{\micro g} de \omterm{def:b3:genetic-engineering:tools}{plasmide}.
\textbf{6.} Le clonage n'exige que la réplication, assurée par
l'origine du \omterm{def:b3:genetic-engineering:tools}{plasmide}~; le promoteur humain n'est pas lu par l'ARN
polymérase bactérienne, de sorte que pour l'expression il faut placer
devant la séquence codante (sans intron) un promoteur bactérien et un
site de fixation du ribosome.
\textbf{7.} $N_{0} = N_{T}/2^{C_{T}}$~: $10^{12}/2^{28} \approx 3700$
copies~; à $C_{T} = 35$, $10^{12}/2^{35} \approx 29$ copies.
\textbf{8.} $2^{n} = 10^{12}$~: $n = 39.9$, quarante cycles.
\textbf{9.} $10^{12}/2^{38} \approx 3.6$ copies. Au-delà d'environ
$40$ cycles, une seule molécule contaminante, ou des artefacts
d'\omterm{def:b3:genetic-engineering:pcr}{amorces}, atteignent le seuil, et un échantillon de moins d'une copie
n'a aucun sens.
\textbf{10.} Une molécule franchit le seuil au cycle $40$~: un positif.
On l'empêche par des salles et du matériel séparés pour la préparation
et pour la manipulation des produits, un flux à sens unique, des
témoins négatifs dans chaque course, et la destruction enzymatique des
amplicons reportés.
\textbf{11.} $3700/0.4 \approx 9000$ copies d'ARN.
\textbf{12.} Rapporté $2^{3.3} \approx 10$ fois~; vrai $1.9^{3.3} =
e^{3.3\ln 1.9} \approx 8.3$ fois.
\textbf{13.} $40\times \qty{35}{\micro g} = \qty{1.4}{mg}$ par jour,
\qty{0.51}{g} par an~; $\times 10^{7}$~: \qty{5100}{kg} par an.
\textbf{14.} $5.1\times 10^{6}/5800 \approx 880$ mol~; $5.3\times 10^{26}$
molécules.
\textbf{15.} $5.1\times 10^{6}$ g $/ (20$ g/L$) = 2.6\times 10^{5}$ L
$= \qty{260}{m^3}$ --- le dixième de la piscine.
\textbf{16.} $5100\times 8000 = 4.1\times 10^{7}$ kg de pancréas, à
\qty{0.1}{kg} l'un~: $4\times 10^{8}$ porcs par an.
\textbf{17.} L'hormone recombinante est identique à l'humaine, donc
elle ne suscite pas d'anticorps et ne cause pas de réaction
allergique~; elle ne véhicule aucun pathogène animal (virus, prions)~;
l'approvisionnement ne dépend pas de l'abattage~; et la séquence peut
être améliorée.
\textbf{18.} Les résidus qui touchent le récepteur sont inchangés, donc
la fixation et la signalisation sont celles de l'insuline~; les résidus
modifiés se trouvent à la surface où les molécules d'insuline
s'associent en dimères et en hexamères, et les changer change la
vitesse à laquelle le dépôt injecté se défait en monomères
absorbables.
\textbf{19.} $6.4\times 10^{9}\times 4^{-20}/16 \approx 3.6\times
10^{-4}$~: aucune correspondance exacte fortuite.
\textbf{20.} $1 + 60 + 1710 + 27\times 1140 = \num{32551}$ séquences~;
$\num{32551}\times 3.6\times 10^{-4} \approx 12$ \omterm{prop:b3:genetic-engineering:editors}{sites hors cible}
fortuits.
\textbf{21.} $10^{7}\times 10^{-4} = 1000$ cellules. Une cellule souche
vit et se divise pendant toute la vie du patient~; l'une d'elles qui a
perdu un suppresseur de tumeur, ou acquis une translocation, peut
fonder un clone qui devient une leucémie.
\textbf{22.} $p_{1} = 0.0019$, $p_{2} = 0.0036$, $p_{3} = 0.0069$~; une
croissance géométrique de $1.9$ par génération donne $\ln 500/\ln 1.9
\approx 10$~; en itérant la récurrence, la fréquence dépasse $0.5$ à la
génération $11$ ($0.45 \to 0.68$).
\textbf{23.} Environ un an. Un coût de \qty{20}{\%} chez les
homozygotes est inférieur à $c = 0.9$, de sorte que le forçage se fixe
tout de même, plus lentement, tandis que la valeur sélective moyenne de
la population baisse --- ce qui est le but d'un forçage de
suppression~; un coût supérieur à $c$ le bloquerait à une fréquence
intermédiaire.
\textbf{24.} $10^{6}\times 0.1 = 10^{5}$ allèles de résistance en une
génération. Incoupables et, s'ils ne coûtent rien là où le forçage
coûte, ils sont favorisés et se propagent, et le forçage est éliminé
--- d'où le choix de cibles où les produits de jonction d'extrémités
sont létaux, avec plusieurs guides à la fois.
\textbf{25.} Environ $3700$ copies dans le premier échantillon~; quelque
\qty{260}{m^3} de fermenteur par an pour l'insuline du monde~; environ
$11$ générations, soit une année de moustique, pour le forçage.
\end{solution}
